\clearpage
\section{Podstawy teoretyczne i przegląd literatury}
\subsection{Sztuczne sieci neuronowe}
Jedną z najprostszych definicji sztucznych sieci neuronowych sformułował dr Robert Hecht-Nielsen, jeden z ojców uczenia maszynowego. Zgodnie z jego definicją sztuczna sieć neuronowa to system obliczeniowy złożony z wielu prostych, silnie połączonych ze sobą elementów, które przetwarzają informacje poprzez dynamiczną reakcję na wejścia z zewnątrz\cite{10.5555/38292.38295}. Można to pojęcie sprecyzować i przedefiniować w oparciu o współczesną realizację sieci neuronowych --- sztuczna sieć neuronowa to zbiór jednostek przetwarzających, wykonujących przekształcenia przeważnie nieliniowe, czasami liniowe. Jednostki te nazywa się neuronami, ponieważ występuje analogia między nimi a komórkami nerwowymi organizmu, czyli przetwarzanie wielu wejściowych informacji w jedną informację, podobnie jak neuron odbiera wiele bodźców przez dendryty i produkuje na ich podstawie aktywację do aksonu. Jednostka sieci neuronowej robi to poprzez ważenie i sumowanie tych informacji oraz, przeważnie, poddaniu wyniku nieliniowej transformacji. Wynika z tego, że sieć neuronowa jest w istocie złożoną funkcją, będącą kompozycją wielu mniejszych funkcji, których współczynniki są dostosowywane (uczone) w procesie treningu

Samodzielne neurony nie zapewniają zdolności do modelowania złożonych wzorców, dlatego grupuje się je w struktury, tworzące sieć. Najbardziej podstawową jest warstwa neuronów, czyli grupa równoległych neuronów na tym samym poziomie głębokości, nie połączonych między sobą, które przetwarzają dane wejściowe, lub aktywacje z poprzedniego poziomu głębokości, a ich aktywacje zasilają kolejne warstwy sieci, lub są jej danymi wyjściowymi. Kompozycja warstw, często różnych, uzupełniona o warstwy normalizujące i warstwy aktywacji tworzą moduły, nazywane również blokami. Na jeszcze wyższym poziomie wymienione dotychczas struktury łączą się w całe architektury. Architektura sieci neuronowej to plan połączenia i organizacji wszystkich warstw oraz modułów w sieci, przepływu danych przez jej elementy, liczbę parametrów, czy zastosowanie technik, które usprawniają trening oraz działanie sieci. Na potrzeby niniejszej pracy w kolejnych podrozdziałach zostaną wprowadzone wybrane struktury, po które sięgnięto opracowując rozwiązanie będące tematem tej pracy.

\subsubsection{Warstwa w pełni połączona}
Warstwa w pełni połączona, jak sama nazwa wskazuje, odbiera aktywacje od wszystkich neuronów poprzedzającej warstwy (lub danych wejściowych sieci), waży je i sumuje, dodając stałą przesunięcia. Są proste i mają dużą moc aproksymacyjną, ale łatwo się przeuczają, dlatego współcześnie rzadko używa się ich samodzielnie --- zwykle są elementem bardziej złożonej struktury lub ostatnią warstwą sieci oraz są uczone z regularyzacją przez wyłączanie losowych neuronów (ang. dropout).

\subsubsection{Perceptron wielowarstwowy}
Z co najmniej dwóch trenowalnych, w pełni połączonych warstw neuronów, o nieliniowych aktywacjach, można zbudować perceptron wielowarstwowy. To prosty, ale potężny model, który --- zgodnie z twierdzeniem o uniwersalnej aproksymacji --- może z określoną precyzją aproksymować dowolną funkcję ciągłą w ograniczonej dziedzinie. Ta własność pozwoliła perceptronowi wielowarstwowemu (ang. Multilayer Perceptron, MLP) stać się podstawą wielu modeli do klasyfikacji, regresji, przetwarzania języka naturalnego, zadań w wizji komputerowej i wielu innych.

\subsubsection{Warstwa splotowa}
Stosowana od dekad w analizie obrazów operacja splotu okazała się niezwykle potężna wraz z rozwojem sieci neuronowych. W klasycznych implementacjach ręcznie projektowane filtry umożliwiały wygładzanie i wyostrzanie obrazów, detekcję krawędzi, a w tych bardziej złożonych segmentację obrazów i ekstrakcję cech. Neuronowe warstwy splotowe idą o krok dalej --- wykorzystują od dziesiątek, do tysięcy filtrów, strojonych w ramach treningu sieci neuronowej, które pozwalają na o wiele dalej idącą analizę. Splotowe siec neuronowe, w szczególności tak zwane głębokie, czyli mające od kilkunastu do kilkudziesięciu warstw, są uczone identyfikowania złożonych kształtów, tekstur, wzorców przestrzennych i łączyć je w wysokopoziomową umożliwiającą interpretację złożonych obiektów.

Technicznie różnią się istotnie od warstw w pełni połączonych. W skład warstwy wchodzi wiele kanałów, a każdy z nich korzysta z jednego, wyuczonego filtra. Wszystkie neurony w kanale współdzielą ten filtr i operują na danych trójwymiarowych, ale, co szczególnie odróżnia je od neuronów w warstwach w pełni połączonych, obliczają aktywacje jedynie na wycinku wejściowego tensora, zwanym polem recepcyjnym. W praktyce jeden neuron przyjmuje wejście o szerokości i wysokości filtra, na głębokości wszystkich kanałów jednocześnie. Aktywacja jednego neuronu tworzy jeden ,,piksel`` w ,,obrazie`` wyjściowym. Słowa piksel i obraz są ujęte w cudzysłowie, ponieważ w istocie warstwy splotowe produkują obiekty nieco inne niż obrazy, czyli mapy cech. Podczas gdy cyfrowe obrazy są tensorem o jakiejś szerokości i wysokości oraz (przeważnie) trzech kanałach głębi, np. odpowiadającym czerwieni, zieleni i błękitu w obrazie, to mapy cech są mniejsze przestrzennie, ale znacznie głębsze. Głębokość ta odpowiada liczbie filtrów w warstwie, więc jak wspomniano wcześniej, może sięgać tysięcy kanałów.

\subsubsection{Autoenkoder i autoenkoder wariacyjny}
Kiedy zadaniem jest nauczenie sieci neuronowej na podstawie danych, które nie są etykietowane, nazywa się to uczeniem nienadzorowanym. W istocie celem tego procesu jest stworzenie modelu danych, który odkryje ukryte struktury, wzorce i zależności. Tak nauczony model może następnie zostać wykorzystany do klastrowania danych, detekcji anomalii czy redukcji wymiarowości danych. Do takich zadań została opracowana architektura autoenkodera, czyli nienadzorowanej sieci zawierającej funkcję kompresującą --- koder --- i rekonstruującą --- dekoder.

Definicja autoenkodera nie stawia ograniczeń wobec typów danych, na jakich może pracować, pod warunkiem, że zostaną przekształcone do postaci liczbowej. Stąd zarówno koder jak i dekoder mogą się składać z warstw w pełni połączonych, splotowych, rekurencyjnych, czy większych bloków jak perceptron wielowarstwowy. Cechą wspólną każdego autoenkodera jest fakt, że stopniowo kompresują dane wejściowe, najczęściej poprzez zmniejszanie wymiarów warstw --- autoenkoder niedopełniony --- lub karania sieci za każdy użyty neuron --- autoenkoder rzadki.

Szczególnym typem autoenkoderów są autoenkodery wariacyjne. W klasycznym podejściu model w mapuje dane wejściowe na jakiś punkt w przestrzeni ukrytej. To sprawia, że przestrzeń ukryta nie jest ciągła, co niesie za sobą dwa główne skutki. Punkty leżące blisko w przestrzeni ukrytej mogą mapować się na zgoła inne zestawy danych, a pobranie losowej próbki przeważnie nie mapuje się na nic, poza szumem. Odpowiedzią na to jest mapowanie danych wejściowych na jakąś próbkę rozkładu, którego parametry są wyznaczane w trakcie treningu. To podejście pozwala na pobieranie z reprezentacji ukrytej próbek, których nawet model nie widział, czyli generowanie nowych danych.

\subsubsection{Transformer}
Szczególną popularność w drugiej dekadzie XXI wieku zyskała architektura transformera z kluczowym dla jej istnienia mechanizmem atencji\cite{attentionisallyouneed2017}. Powstała w celu poprawy efektywności przetwarzania danych sekwencyjnych, ponieważ rozwiązania znane do tamtego czasu pozwalały jedynie na przetwarzanie elementów w kolejności, a mechanizm atencji przetwarza duże części sekwencji jednocześnie. Nazwa architektury wywodzi się stąd, że pierwotnie ich zastosowaniem było transformowanie sekwencji wejściowej w sekwencję wyjściową, na przykład w tłumaczeniach. Obecnie transformery wykorzystuje się do różnych zadań na tekście i to one stały się podstawą dla wielkich modeli językowych, choć sprawdzają się także w pracy z sygnałami dźwiękowymi i danymi wizualnymi.

Działanie sieci zaczyna się od podziału sekwencji na mniejsze fragmenty --- tokeny --- i przekształceniu ich w wektory liczbowe, tzw. wektory cech (ang. embeddings), czyli liczbową reprezentację tokenów, która może być przetwarzana przez sieć. Ponieważ cechą transformerów jest to, że nie tylko uczą się znaczenia pojedynczych tokenów oraz ich sekwencji, a także z faktu, że nie przetwarzają ich sekwencyjnie, a wsadowo, to znaczenie tokenów jest różnicowane względem pozycji w sekwencji, za pomocą kodowania pozycyjnego. Dodatkowy wektor, wyliczany na podstawie funkcji sinusoidalnych o różnych częstotliwościach, jest dodawany do wektora cech.

Tak przygotowane wektory trafiają do serca sieci typu transformer, czyli mechanizmu atencji. Tam każdy element sekwencji jest analizowany pod kątem powiązania z wszystkimi pozostałymi elementami, także z samym sobą. Dzieje się to przez transformacje wektorów elementów sekwencji do trzech nowych wektorów dla każdego z elementów:
\begin{itemize}
	\item \textbf{wektor zapytania (Q)} --- można go interpretować jako reprezentację zapytania, jakie ten wektor złoży kluczom innych elementów sekwencji
	\item \textbf{wektor klucza (K)} --- wektor komplementarny dla wektora zapytania
	\item \textbf{wektor wartości (V)} --- ten wektor niesie wartość elementu
\end{itemize}
Wektory Q i K są poddawane iloczynowi skalarnemu, na którego podstawie oblicza się macierz prawdopodobieństwa zależności pomiędzy dwoma elementami. Następnie wyznaczone prawdopodobieństwa są mnożone z wektorem wartości dając wynik mechanizmu atencji. Wektory Q i K są tworzone z wejściowej sekwencji poprzez przemnożenie je przez dedykowane, wytrenowane macierze.

W zależności od wag w macierzach, mechanizm atencji może w różny sposób znajdować zależności, dlatego stosuje się tak zwany wielogłowicowy mechanizm atencji. Sieć oblicza uwagę w wielu równoległych głowicach, a następnie łączy wyniki. W kolejnym kroku wynik mechanizmu atencji poddaje się normalizacji, przepuszcza przez MLP, który przetwarza nieliniowe zależności i znów normalizuje. Część z mechanizmem atencji, perceptronem i normalizacją tworzy spójny blok, który w oryginalnej architekturze jest powtórzony kilkukrotnie w sposób szeregowy. Łącząc wszystkie opisane elementy otrzymujemy koder transformera, czyli komponent, którego celem jest zakodowanie sekwencji wejściowej w ukrytą reprezentację.

Drugą częścią oryginalnej architektury transformera jest dekoder. W przeciwieństwie do autoenkodera nie przyjmuje on na wejściu bezpośrednio reprezentacji z kodera, a inną sekwencję. W praktyce, zawsze jest to fragment sekwencji wyjściowej, który transformer już utworzył w procesie przekształcania sekwencji wejściowej. Wobec tego, w wytrenowanym modelu w trybie wnioskowania na początku procesu jest to pusta sekwencja, uzupełniana autoregresyjnie w każdej iteracji dekodera o nowy token. W trakcie treningu podaje się dekoderowi oczekiwany ciąg docelowy, aby możliwe było obliczenie błędu pomiędzy nim, a ciągiem utworzonym przez dekoder.

Budowa dekodera jest zbliżona do kodera. Przebieg znów inicjuje utworzenie wektora cech sekwencji oraz kodowanie pozycyjne, choć jedynie dla nowych elementów. Analogicznie jak w koderze, kolejnym krokiem również jest mechanizm atencji, jednak tym razem nieco zmodyfikowany. Skoro dekoder podczas treningu dysponuje całą sekwencją docelową, ale jego zadaniem jest budowanie jej od początku, poprzez dopisywanie właściwych tokenów do z początku pustej sekwencji, to mechanizm atencji może korzystać do analizy tylko z poprzedzających elementów od aktualnie przewidywanego --- dlatego tokeny następujące są maskowane. Ten etap tworzy wewnętrzną reprezentację tego, co już zostało zapisane w sekwencji i jaką informację niesie ten ciąg, czyli wzbogaca już powstały fragment sekwencji o kontekst.

Z tej informacji, która jest reprezentowana przez wektory wzbogacone o informację kontekstową tworzy się znów wektor Q dla kolejnego mechanizmu atencji. Jednak wektory K i V uzyskuje się z reprezentacji utworzonej przez koder. Dlatego kolejny etap nazywa się krzyżowym mechanizmem atencji (ang. cross-attention), ponieważ krzyżuje się w nim informacja z wejściowej sekwencji, która zasiliła koder, z informacją stopniowo odkodowywaną przez dekoder. Celem cross-attention jest zważenie istotności elementów reprezentacji ukrytej wygenerowanej przez koder, aby na podstawie tych wag stworzyć nowy wektor kontekstowy, niosący informację z dekodera, wzbogaconą o informację z kodera, która jest następnie wykorzystywane do wygenerowania następnego tokenu w sekwencji dekodera.

\begin{sloppypar}
Blok, w którym znajdują się mechanizmy atencji oraz warstwy normalizacyjne, za\-koń\-czo\-ny jest MLP i powtórzony kilku-kilkunastu krotnie dla pogłębienia analizy. Z przetworzonej sekwencji wektorów wyciągany jest ostatni z nich i poddaje się go transformacji liniowej z rozmiaru modelu, czyli długości wektora cech, do wymiaru słownika modelu i klasyfikacji najbardziej prawdopodobnego następnego tokenu ze słownika. Predykowany token jest dopisywany do sekwencji i proces się powtarza, do momentu, gdy zostanie dopisany token \texttt{<end>} kończący generowanie.
\end{sloppypar}

\subsubsection{Transformer wizyjny z przesuwnym oknem}
Jako ostatnią wartą omówienia w niniejszej pracy należy wskazać architekturę będącą jej podstawą. Opisany w poprzedniej sekcji transformer został zaprojektowany do pracy nad tekstem, jednak szybko pojawił się nurt w badaniach, aby wykorzystać tę strukturę do zadań w wizji komputerowej. Pierwszy duży sukces odniósł model transformera wizyjnego (ang. Vision Transformer, ViT), który był w stanie pokonać najlepsze modele splotowe w teście wydajności w zadaniu klasyfikacji. Był tylko jeden problem --- ViT wymaga treningu przy użyciu ogromnej mocy obliczeniowej i setek milionów obrazów, aby osiągnąć tak dobre efekty. Dlatego też w niniejszej pracy z niego zrezygnowano, na rzecz wariacji inicjalnego transformera wizyjnego, czyli transformera z oknem przesuwnym (ang. Shifted Window Transformer, Swin Transformer).

Zasada działania architektury Swin Transformera jest zbliżona do oryginalnego transformera, ale ma zasadnicze różnice. Pierwszą polega na tym, że podczas gdy klasyczny transformer dzieli tekst na tokeny, które to są następnie odwzorowywane na odpowiadające im wektory cech. W przypadku transformera wizyjnego tokenizacja odbywa się najczęściej za pomocą warstwy splotowej, która przekształca obszar obrazu (w przypadku Swin Transformera są to łatki o wymiarze \(4 \times 4 \) pikseli) w jednowymiarowy wektor o długości, na przykład, 96 wartości. Funkcja ta pełni jednocześnie rolę wektora cech, z pominięciem kodowania pozycyjnego. Efektem tej operacji jest przekształcenie obrazu, o najczęściej trzech kanałach RGB, do trójwymiarowej reprezentacji o głębokości, równej długości tych wektorów. W tym miejscu łatwo dostrzec analogię do mapy cech w sieciach splotowych. 

Uzyskana siatka wektorów trafia do bloku z mechanizmem atencji. Tu leży największa różnica pomiędzy ViT a Swin Transformerem. Podczas gdy ten pierwszy ,,karmi'' mechanizm atencji całą siatką łatek, ten drugi dzieli ją na mniejsze, niezachodzące na siebie okna, np. 7x7 łatek, dzięki czemu drastycznie maleje złożoność obliczeniowa. Efektem niepożądanym tego zabiegu jest obniżenie zdolności do modelowania dalekodystansowych zależności. Odpowiedzią twórców na ten problem jest mechanizm przesuwnego okna. W kolejnym mechanizmie atencji siatka okien zostaje przesunięta względem siatki łat o połowę rozmiaru okna, co umożliwia łączenie informacji między sąsiednimi oknami, poprawiając modelowanie zależności dalekodystansowych. Zmodyfikowany mechanizm atencji wprowadza również zależności pozycyjne między łatkami. Kodowanie pozycyjne z początku sieci zostało zastąpione dodaniem uzupełniającej informacji do łatek o relatywnym położeniu względem innych łatek w oknie. Sekwencja mechanizm atencji -> MLP -> przesunięty mechanizm atencji -> MLP nazywa się blokiem transformera Swin.

Swin Transformer składa się z czterech większych bloków, zwanych etapami. Pierwszy etap to utworzenie wektora cech i dwa bloki transformera. W kolejnych trzech etapach w pierwszym kroku dokonuje się łączenie łatek, w wyniku czego każdy z dwóch wymiarów siatki zmniejsza się dwukrotnie, a jej głębokość dwukrotnie zwiększa. To drugi z zabiegów niwelujących negatywny efekt wprowadzenia okien. Poprzez łączenie łatek pojedyncze okno obejmuje obszar z wielu okien z poprzedniego etapu, w efekcie czego informacje o zależnościach z odległych okien mogą być propagowane. W następnej kolejności nowa siatka okien i łatek jest przetwarzana w kolejnych blokach danego etapu transformera. Wynikiem pracy każdego z etapów transformera Swin jest hierarchiczna reprezentacja obrazu, w której wzorce lokalne i globalne są coraz lepiej modelowane. Finalna reprezentacja jest gotowa do dalszych zadań, takich jak klasyfikacja, segmentacja czy detekcja obiektów, a w przypadku niniejszej pracy, dalszym operacjom w kierunku kompresji.

\subsubsection{Funkcje aktywacji}

\subsection{Kompresja danych}
Kompresję danych można zdefiniować jako zmniejszenie ich rozmiaru, czyli liczby bitów potrzebnych do ich zakodowania i przeniesienia informacji którą reprezentują. Podstawowym wymaganiem jest, aby istniała możliwość wykonania na skompresowanej reprezentacji operacji odwrotnej, nazywanej dekompresją, która pozwala na odtworzenie informacji. Nie oznacza to jednak, że kompresja zawsze jest odwracalna. Terminologia tej dziedziny wyróżnia bowiem zarówno kompresję odwracalną (bezstratną), jak i nieodwracalną (stratną). W przypadku pierwszej możliwe jest odtworzenie informacji w procesie dekompresji dokładnie takiej, jaka została poddana kompresji. a w drugim przypadku dopuszczamy możliwość utraty części informacji.

Zastosowanie kompresji bezstratnej jest konieczne wszędzie tam, gdzie niezachowanie pełnej informacji wiązałoby się z uszkodzeniem jej lub wprowadzeniem błędu, który jest wysoce niepożądany z powodu krytycznego znaczenia informacji. Przykładów jest tutaj wiele, a jednym z najczęstszych, z jakimi spotykają się użytkownicy cyfrowych technologii (pomijając obrazy, które zostaną znacznie szerzej opisane w dalszej części pracy), są archiwa, np. ZIP. Użytkownik umieszczający pliki w archiwum chce zmniejszyć ich rozmiar, ale oczekuje, że po rozpakowaniu archiwum otrzyma pliki o niezmienionej zawartości i funkcjonalności. Gdyby było inaczej i w archiwum umieszczony zostałby np. program, to utrata informacji spowodowałaby jego uszkodzenie. i w krytycznych zastosowaniach, nawet jeżeli wprowadzenie błędu nie spowoduje uszkodzenia, to może nieść poważne negatywne skutki. Niech przykładem będzie scenariusz kompresji danych medycznych --- utrata precyzji w zapisie tej informacji może mieć wpływ na nieprawidłową diagnozę, przeoczenie schorzenia. 

Argumentem na korzyść idei kompresji stratnej jest możliwość znacznie silniejszej redukcji rozmiaru danych, przy jednoczesnym zachowaniu użyteczności informacji. Kluczem jest znalezienie odpowiedniego kompromisu w zadanym kontekście --- przykładowo w przypadku przesyłania strumieniowego uzasadnione jest przesunięcie tego kompromisu na rzecz wyższej przepływności kosztem jakości, zaś kiedy mowa o przechowywaniu informacji w pamięci trwałej, która nie będzie przesyłana wąskopasmowym łączem sieciowym, można skupić się na jakości. To nieodwracalne podejście doskonale nadaje się do użycia dla danych multimedialnych. Ludzka percepcja jest daleka od perfekcji --- bodźce wzrokowe czy dźwiękowe z jednej strony są odbierane indywidualnie przez każdego człowieka, z drugiej wszyscy mamy te same ograniczenia wynikające z fizjologicznego uwarunkowania. Dla bodźców dźwiękowych będzie to ograniczony zakres słyszalności w dziedzinie częstotliwości, niezdolność do rozróżnienia dźwięków o bardzo zbliżonej częstotliwości, czy możliwość pominięcia bardzo krótkich dźwięków. W przypadku treści wizualnych reguły są bardzo podobne. Punktowe różnice mogą zostać pominięte przez ludzkie oko, barwy o niewielkiej różnicy są słabo rozróżniane (bardziej dostrzegalne są różnice w jasności), a uwaga odbiorcy jest wybiórcza – bardziej skupiona na kształtach, postaciach, obiektach pierwszoplanowych niż na teksturach i tle. Tę wiedza pozwala na redukcję wierności odwzorowania mniej istotnych elementów treści multimedialnych, dzięki czemu możliwe jest ich silniejsze skompresowanie, przy zachowaniu użyteczności i zadowalającej percepcji u odbiorcy.

\subsection{Klasyczne metody kompresji}
\subsubsection{JPEG}
Wyżej wymienione cechy stały się fundamentem dla kodeku JPEG. Kluczem do sukcesu tej metody jest wykorzystanie faktu, że ta część informacji na obrazie, która jest najważniejsze dla ludzkiego oka jest zakodowana w niskich częstotliwościach, a wysokie częstotliwości odpowiadają drobnym szczegółom, które są mniej istotne dla percepcji. Opracowana w roku 1972 dyskretna transformacja kosinusowa (ang. discrete cosinus transform, DCT) pozwoliła na analizę obrazu w dziedzinie częstotliwości i wykorzystanie wspomnianej właściwości. W kolejnych latach pracowano nad usprawnieniem metody i jej efektywną aplikacją. Powołana w 1986 roku grupa naukowców z Joint Photographic Experts Group opracowała skuteczną metodę wykorzystującą m. in. DCT, ale również technikę podpróbkowania oraz kodowanie entropijne. Wprowadzony w 1992 roku kodek stratnej kompresji obrazów po dziś dzień znajduje szerokie zastosowanie w transmisji cyfrowych obrazów. Za jego popularnością stoi też prostota algorytmu. Proces kompresji z użyciem kodeka JPEG składa się z kilku etapów.

Pierwszym krokiem jest przekształcenie obrazu z przestrzeni barw RGB (czerwony, zielony, niebieski, ang. \textbf{R}ed, \textbf{G}reen, \textbf{B}lue) do przestrzeni YCbCr. W tej reprezentacji składowa Y odpowiada za jasność (luminancję) obrazu, natomiast komponenty Cb i Cr zawierają informacje o barwie (chrominancji), czyli odpowiednio o odchyleniu koloru niebieskiego i czerwonego. Taka transformacja pozwala na zdekorelowanie informacji o jasności i barwie. W kolejnym podetapie, dokonuje się podpróbkowania składowych barwowych – najczęściej zmniejsza się ich rozdzielczość przestrzenną (np. poprzez zastosowanie podpróbkowania w stosunku 4:2:0 lub 4:2:2), co prowadzi do znacznego zmniejszenia ilości danych bez zauważalnego pogorszenia jakości wizualnej obrazu. Dzięki tej technice możliwe jest wykorzystanie w praktyce właściwości ludzkiego systemu wzrokowego, że różnice tonalne są dostrzegane słabiej od różnic w jasności.

Następnie przetworzony obraz dzieli się na mniejsze bloki –--- standardowo o rozmiarze \(8\times8\) pikseli. Każdy z tych bloków jest analizowany z użyciem dyskretnej transformaty DCT. Transformata ta przekształca dane z dziedziny przestrzennej, czyli jasności lub barwy poszczególnych pikseli, do dziedziny częstotliwości. Rezultatem tej operacji jest zestaw współczynników, z których większość odpowiada za szczegóły i drobne zmiany w obrazie, natomiast tylko kilka niesie główną informację o jego strukturze. DCT umożliwia więc reprezentację obrazu w sposób, który jest bardziej podatny na skuteczną kompresję.

Kolejnym etapem jest kwantyzacja otrzymanych współczynników. W tym procesie wartości wynikające z transformacji są zaokrąglane do najbliższych określonych poziomów, co skutkuje utratą pewnej ilości informacji --- dlatego kompresja JPEG jest klasyfikowana jako metoda stratna. Jednak dzięki temu, że wyższe częstotliwości (czyli szczegóły) mają mniejsze znaczenie percepcyjne, można je silniej kwantować. Umożliwia to znaczącą redukcję rozmiaru danych przy ograniczonym wpływie na postrzeganą jakość obrazu.

Ostatnim etapem jest bezstratna kompresja skwantowanych danych. Dzięki niej możliwe jest uzyskanie jeszcze większego stopnia kompresji, ale bez dodatkowej utraty informacji. Do tego celu wykorzystuje się kodowanie entropijne, takie jak np. kodowanie Huffmana lub kodowanie arytmetyczne. Wykorzystując informację o tym, jak często występują poszczególne wartości, można zakodować częstsze z nich za pomocą krótszych ciągów, a rzadsze --- dłuższymi.

Proces dekompresji obrazu JPEG polega na odtworzeniu możliwie jak najbardziej zbliżonego do oryginalnego, na podstawie danych skompresowanych podczas kodowania. Rozpoczyna się on od odczytania i zdekodowania skompresowanych danych za pomocą odwrotnego kodowania entropijnego (np. dekodera Huffmana), co pozwala odzyskać skwantowane współczynniki DCT. Następnie współczynniki te są poddawane odwrotnej kwantyzacji. Należy zauważyć, że informacji utraconych podczas wcześniejszego zaokrąglania nie da się odzyskać --- co stanowi istotę stratnej natury tego kodeka. Kolejnym krokiem jest zastosowanie odwrotnej transformaty DCT (Inverse Discrete Cosine Transform, IDCT), która przekształca dane z powrotem z dziedziny częstotliwości do dziedziny przestrzennej, umożliwiając odtworzenie rozkładu jasności i barw w blokach \(8 \times 8 \) pikseli. Dalej, jeśli zastosowano podpróbkowanie chrominancji, wartości w kanałach Cb i Cr są interpolowane do pełnej rozdzielczości i razem z luminancją Y konwertowane z przestrzeni YCbCr z powrotem do przestrzeni RGB. W wyniku tych operacji uzyskujemy końcowy obraz, który choć nie jest identyczny z oryginałem, zachowuje znaczną część inicjalnej informacji i ma zadowalającą jakość wizualną.

Chociaż JPEG ma wiele zalet jak wydajność, prostota i jest powszechnie stosowanym formatem kompresji obrazów, posiada również swoje ograniczenia:
również swoje ograniczenia:

\begin{itemize}
	\item Artefakty blokowe --- ponieważ obraz jest dzielony na bloki \(8 \times 8 \) pikseli, przy dużym stopniu kompresji mogą pojawić się charakterystyczne granice między blokami.
	
	\item Brak wsparcia dla przezroczystości --- JPEG nie obsługuje kanału alfa, co ogranicza jego zastosowanie np. w grafice internetowej wymagającej przezroczystego tła.
	
	\item Słaba wydajność w przypadku grafik z ostrymi krawędziami --- kompresja JPEG gorzej radzi sobie z obrazami zawierającymi tekst, wykresy czy grafikę wektorową, gdzie dominują ostre przejścia tonalne.
	
	\item Mimo tych wad, JPEG pozostaje jednym z najważniejszych standardów w dziedzinie kompresji obrazów rastrowych --- głównie dzięki znakomitemu kompromisowi pomiędzy stopniem kompresji a jakością wizualną w przypadku fotografii, będąc jednocześnie bardzo lekkim.
\end{itemize}

\subsubsection{JPEG 2000}
Po opublikowaniu kodeka grupa JPEG podjęła działania w celu stworzenia bezstratnego odpowiednika --- JPEG LS. Swój wkład postanowili dołożyć badacze z japońskiej firmy Ricoh przysyłając zespołowi standaryzującemu autorski algorytm CREW (\textbf{C}ompression with \textbf{R}eversible \textbf{E}mbedded \textbf{W}avelets). Choć w pracach nad standardem bezstratnym to rozwiązanie zostało odrzucone na rzecz algorytmu LOCO-I (\textbf{Lo}w \textbf{Co}mplexity \textbf{Lo}ssless \textbf{Co}mpression for \textbf{I}mages), dostrzeżono jego potencjał i uruchomiono prace nad nowym kodekiem w 1997 roku. Przyjęto szereg warunków jakie miał spełniać\cite{Taubman2002}:
\begin{enumerate}
	\item \label{jakosc-bitrate} Wysoka jakość w niskim budżecie bitowym.
	\item \label{glebia} Obsługa obrazów w szerokim zakresie głębi bitowej (od 1 do 16 bitów na składową piksela).
	\item \label{progresja} Progresywne strumieniowanie względem jakości i rozdzielczości obrazu.
	\item \label{stratna-bezstratna} Jednolita obsługa kompresji stratnej i bezstratnej w ramach jednego formatu.
	\item \label{swobodny-dostep} Dostęp do dowolnego obszaru obrazu ze strumienia bitowego bez dekodowania całości.
	\item \label{odpornosc} Odporność na błędy bitowe w transmisji.
\end{enumerate}
JPEG 2000 spełnia te warunki za sprawą przemyślanych technik w nim zastosowanych. Podstawą algorytmu jest zastąpienie transformacji kosinusowej transformacją falkową --- DWT (Discrete Wavelet Transform). Falka, w przeciwieństwie do fal sinusowych w DCT, oscyluje w skończonym przedziale czasowy. Znacznie lepiej nadaje się do analizowania lokalnych zmian w obrazach od funkcji sinusowych, które oscylują w każdym punkcie dziedziny czasu. Kodek JPEG 2000 wykorzystuje dwa rodzaje:
\begin{itemize}
	\item \textbf{Falka 9/7 Daubechies} --- 9 wag filtru dolnoprzepustowego i 7 górnoprzepustowego, które są liczbami rzeczywistymi, co otwiera drogę do redukowania precyzji w kwantzacji --- wykorzystywana w trybie stratnym.
	\item \textbf{Falka 5/3 LeGall} --- 5 wag filtru dolnoprzepustowego i 3 górnoprzepustowego, które są liczbami całkowitymi --- wykorzystywana w trybie bezstratnym.
\end{itemize}

Oryginalny obraz, przed poddaniem DWT, jest rozkładany na składowe --- każdy kanał jest poddawany analizie z osobna --- a każda ze składowych dzielona na części, tzw. kafelki. Standard oferuje w tej materii dużą dowolność, bowiem rozmiar kafelka nie jest narzucony z góry i w praktyce jeden kafelek może obejmować rozmiar całego obrazu. Pojedynczą jednostkę kodowania poddaje się filtrowaniu w sześciu krokach:
\begin{enumerate}
	\item Filtrem dolnoprzepustowym wzdłuż wierszy, odrzucając co drumgą kolumnę --- otrzymując macierz współczynników L.
	\item Filtrem górnoprzepustowym wzdłuż wierszy, odrzucając co drugą kolumnę --- otrzymując macierz współczynników H.
	\item Macierz L filtrem dolnoprzepustowym wzdłuż kolumn, odrzucając co drugi wiersz --- otrzymując macierz współczynników LL.
	\item Macierz L filtrem górnoprzepustowym wzdłuż kolumn, odrzucając co drugi wiersz --- otrzymując macierz współczynników LH.
	\item Macierz H filtrem dolnoprzepustowym wzdłuż kolumn, odrzucając co drugi wiersz --- otrzymując macierz współczynników HL.
	\item Macierz H filtrem górnoprzepustowym wzdłuż kolumn, odrzucając co drugi wiersz --- otrzymując macierz współczynników HH.
\end{enumerate}
W efekcie powstają cztery macierze o dwukrotnie zmniejszonym wymiarze przestrzennym, z których LL zawiera ogólną, wygładzoną informację o obrazie, LH i HL dodają jej szczegółów, a informacja z HH to informacja o wysokiej częstotliwości --- detale, krawędzie, złożone tekstury (LH wyodrębnia krawędzie poziome, HL krawędzie pionowe, a HH krawędzie ukośne). Innymi słowy dzięki DWT otrzymuje się podział obrazu na pasma o rosnącej częstotliwości. Dodatkowo transformacja falkowa w tym miejscu się nie kończy --- macierz LL poddaje się kolejnej transformacji, dzieląc ją na kolejne stopnie szczegółowości i umożliwiając kolejną iterację. Po kilku obiegach uzyskuje się hierarchiczną strukturę pozwalającą na budowanie obrazu od ogółu do szczegółu.

W następnym kroku współczynniki falkowe w macierzach są kwantowane (z zadanym krokiem kwantyzacji w zależności od poziomu kompresji) i kodowane. W tym miejscu zastosowano drugie kreatywne rozwiązanie pozwalające na elastyczną kontrolę rozmiaru pliku i jakości rekonstrukcji. Zamiast kodować cały zestaw współczynników w podpasmach, JPEG 2000 dzieli je na bloki kodowe. Ich rozmiar może mieć od \(4 \times 4\) do \(1024 \times 1024\) wartości, ale typowo używa się \(32 \times 32\) lub \(64 \times 64\) --- mniejsze bloki pogarszają kompresję, a większe zwiększają złożoność obliczeniową. Kodowanie odbywa się w dwóch fazach. W fazie pierwszej w każdym z bloków wyznacza się najwyższą płaszczyznę bitową, czyli liczbę bitów potrzebną do zapisania największego współczynnika. Wtedy na każdej płaszczyźnie bitowej współczynniki są analizowane jako istotne lub nieistotne --- domyślnie każdy jest nieistotny, dopóki w zapisie tego współczynnika na danej płaszczyźnie nie pojawi się bit 1. Współczynniki istotne i nieistotne na danej płaszczyźnie, mogą być zakodowane taniej, ponieważ statystycznie wartości tej samej istotności sąsiadują ze sobą, więc można je grupować. Bity współczynników na danym planie, wraz z kontekstem sąsiadów, są przesyłane do kodera arytmetycznego, który wykorzystuje te statystyki do efektywnego zakodowania ich w strumieniu bitowym. W zależności od zadanego rozmiaru pliku strumień bitów jest obcinany w optymalnym punkcie w fazie drugiej. Odrzucane są wtedy najmniej znaczące bity, które dają najmniejszą poprawę jakości, w stosunku do kosztu, co efektywnie zmniejsza rozmiar pliku względem precyzji. Technika ta ma w ogólności na celu odrzucenie mniej istotnych dla rekonstrukcji bitów w celu dopasowania się do budżetu bitowego, niezależnie w każdym bloku kodowym, co realizuje założenie nr \ref{jakosc-bitrate}, zwiększa elastyczność kodowania i uodparnia na błędy bitowe w transmisji (izolacja bloków --- izolacja błędów) --- założenie \ref{odpornosc}.

Powyższy opis przedstawia ogólne działanie kodeka z naciskiem techniki kompresji stratnej. Tryb bezstratny jest szczególnym przypadkiem, zawierającym trzy modyfikacje:
\begin{itemize}
	\item Filtr Daubechies 9/7 jest zastąpiony filtrem LeGall 5/3 --- tworzy współczynniki całkowitoliczbowe, które pozwalają na całkowite odwrócenie operacji bez błędów zaokrągleń.
	\item Ponieważ współczynniki są całkowite, a celem jest zachowanie pełnej informacji, kwantyzacja w praktyce nie występuje --- krok kwantyzacji = 1.
	\item Żadne bity ze strumienia kodowego nie są obcinane.
\end{itemize}
Za pomocą tych technik spełniono założenie \ref{stratna-bezstratna}.

Należy dokonać jeszcze dodatkowych wyjaśnień dla pozostałych założeń wspomnianych na początku podrozdziału. Spełnienie punktu nr \ref{glebia} można wyjaśnić krótko. Przedstawiony powyżej opis działania kodeku nigdzie nie przyjmuje sztywnych wartości głębi ani nie narzuca ograniczeń --- transformacja falkowa, kwantyzacja czy kodowanie działają, w dużej mierze, w oderwaniu od liczby bitów. Choć kodek wykorzystuje tę informację, to w tak skrótowym opisie nie ma konieczności się nad tym pochylać.

Uwaga do punktu nr \ref{progresja}. jest następująca: JPEG 2000 realizuje progresje rozdzielczości i jakości poprzez hierarchiczne upakowanie informacji w strumieniu bitowym w postaci pakietów. Ta struktura pozwala na wyodrębnienie poszczególnych podpasm oraz poszczególnych warstw jakości w postaci planów bitowych, umożliwiając zarówno uzyskanie miniatury obrazu, jak i obrazu o pełnej rozdzielczości, ale ograniczonym poziomie odwzorowania detali.

Z kolei dla uzupełnienia punktu \ref{swobodny-dostep}. o swobodnym dostępnie należy dodać, że początek każdego z kafelków jest oznaczony w strumieniu bitowym markerami Start of Tile i Start of Data, co pozwala odkodować fragment obrazu bez odkodowania reszty. Istnieje także w JPEG 2000 mechanizm priorytetyzowania fragmentów obrazu, narzucając im wyższą precyzję kodowania. Realizuje się to za pomocą prostego zabiegu --- przesunięcia bitowego współczynników w bloku kodowym i zapisaniu liczby bitów tegoż przesunięcia w nagłówku obrazu. Zapobiega to obcięciu szczegółów w tym bloku kodowym.


\subsubsection{WebP} \label{webp}

Wraz z rozwojem internetu wąskim gardłem, jeśli chodzi o szybkość ładowania stron okazało się wczytywanie obrazów. Z tego powodu firma Google postanowiła opracować nowy format zapisu zdjęć, który mógłby zastąpić JPEG i przyspieszyć ładowanie stron. Za sukces odpowiada firma On2, od 2010 roku należąca do Google, która od lat 90. pracowała nad kodekami kompresji (choć ich praca skupiona była wokół efektywnego zapisu wideo). W 2010 roku stworzyli format WebM (Web Media) do kompresji audio i wideo, oparty na kodeku VP8 stworzonego przez firmę wcześniej, a techniki, które wykorzystano do efektywnej kompresji pojedynczych klatek, przeniesiono do oddzielnego formatu przeznaczonego dla zdjęć. Cztery miesiące później zaprezentowano format WebP (Web Picture) do kompresji obrazów statycznych i animowanych, w trybach bezstratnym i stratnym. Z uwagi na tematykę pracy pochylę się nad tym drugim. Kodek ten również działa w kilku etapach.

\begin{sloppypar}
Proces kompresji również zaczyna się od konwersji do YCbCr i podpróbkowania, domyślnie w proporcjach 4:2:0. Następnie dokonuje się podziału wejściowego obrazu na bloki \(16\times16\) w kanale luminancji i po \(8\times8\) w kanałach chrominancji.
\end{sloppypar}
	
Szczególnie istotnym elementem procesu jest kolejny krok, w którym dokonywana jest predykcja wartości jasności pikseli na podstawie punktów zakodowanych wcześniej. Predykcje wyznaczane są w oparciu o wartości z bloków powyżej i na lewo od kodowanego i zdefiniowane są różne strategie na to, jak wykorzystać te sąsiadujące bloki do uzyskania nowej wartości. Po wyznaczeniu predykcji koduje się różnicę pomiędzy nią a wartością rzeczywistą. Badania wykazują, że sąsiadujące ze sobą piksele w obrazach naturalnych mają współczynnik korelacji na poziomie 0,9\cite{imagestatistics}, i choć ten współczynnik maleje wraz z odległością, to nadal daje dostateczną przewagę aby predykcja zmniejszyła entropię danych.

Wartości resztkowe są następnie poddawane DCT w blokach \( 4 \times 4 \), a ze względu na ich małą wariancję, wiele współczynników ma wartości bliskie zeru. Poddane w następnym kroku kwantyzacji przez zaokrąglenie, zerują się, co dodatkowo zwiększa efektywność kompresji. Znów, ten krok jest elementarny dla procesu stratnej kompresji, gdyż prowadzi do nieodwracalnej utraty informacji. Skwantowane współczynniki trafiają do kodera arytmetycznego, gdzie są przekształcane w strumień bitów. Ta technika kodowania jest bardziej efektywna niż kodowanie Huffmanna, które jest domyślnym w najpopularniejszych implementacjach JPEG.

Proces dekompresji w formacie WebP podobnie opiera się na założeniu odwrócenia kroków kompresji. Rozpoczyna się od dekodowania arytmetycznego strumienia bitów, które odtwarza skwantowane współczynniki DCT. Te następnie poddawane są dekwantyzacji, która polega na pomnożeniu ich przez te same wartości, przez które były dzielone podczas kwantyzacji. Ten krok przywraca przybliżoną, oryginalną skalę współczynników, jednak bezpowrotnie utracone informacje (wartości wyzerowane podczas kwantyzacji) nie mogą zostać odzyskane. Po dekwantyzacji stosuje się IDCT, która przekształca współczynniki z powrotem z dziedziny częstotliwości do dziedziny przestrzennej, odtwarzając wartości resztkowe dla każdego bloku \(4 \times 4 \). Kluczowym elementem jest odtworzenie finalnych wartości pikseli poprzez dodanie zrekonstruowanych wartości resztkowych do bloków predykcyjnych. Informacje o tym, który tryb predykcji został użyty dla danego bloku, są zapisane w strumieniu danych,
co pozwala dekoderowi na precyzyjne odtworzenie prognozowanych wartości na podstawie już zdekodowanych sąsiednich bloków. Na końcu, po odtworzeniu wszystkich bloków w kanałach luminancji i chrominancji, dane są konwertowane z powrotem z przestrzeni barw YCbCr do RGB, co pozwala na uzyskanie odtworzonego obrazu.

WebP unowocześnia podejście do kompresji obrazów, czerpiąc z zaawansowanych technik opracowanych na potrzeby kodowania wideo. Jego przewaga nad formatem JPEG wynika przede wszystkim z dwóch kluczowych innowacji: predykcji blokowej oraz kodowania arytmetycznego. Predykcja, bazująca na silnej korelacji sąsiadujących pikseli, znacząco redukuje ilość danych, które  muszą zostać zapisane, ponieważ zamiast pełnych wartości kodowane są jedynie niewielkie różnice. Z kolei wydajniejsze kodowanie arytmetyczne pozwala na bardziej zwarte upakowanie tych danych w strumieniu bitów. Połączenie tych metod sprawia, że WebP przy zachowaniu porównywalnej jakości wizualnej oferuje pliki o znacznie mniejszym rozmiarze niż JPEG. To bezpośrednio przekłada się na krótsze czasy ładowania stron internetowych, co było głównym celem jego powstania. Niemniej format ten wciąż ma problemy z kompatybilnością –-- starsze przeglądarki, edytory graficzne i systemy operacyjne mogą nie obsługiwać formatu WebP bez dodatkowych wtyczek lub konwersji.

\subsubsection{BPG}
Dla przeprowadzenia opisu kodeka BPG, należy zacząć od krótkiej historii kodowania plików wideo, która, choć wykracza poza tematykę niniejszej pracy, stanowi część jego własnej historii. Wynika to z faktu, że BPG jest adaptacją standardu kodowania wideo --- HEVC (High Efficiency Video Coding).

Kompresowanie wideo jest niemniej ważnym, a może nawet ważniejszym aspektem transmisji multimediów, stąd od dekad stanowi wyzwanie równoległe dla kompresji obrazów. Od lat 70-tych pracowano nad wieloma rozwiązaniami, jednak w latach 80-tych narodziła się potrzeba standaryzacji, aby urządzenia różnych producentów mogły wymieniać ze sobą informacje. Prym nad pracami wzięły zespoły ekspertów MPEG (Moving Picture Experts Group) --- wydzielona z szerszego aliansu ISO (International Organization for Standarization) oraz IEC (International Electrotechnical Commission) --- oraz VCEG (Video Coding Experts Group) --- wydzielona z ITU-T (International Telecommunication Union Telecommunication Standarization Sector).

Pierwszy istotny na ścieżce rozwoju format H.261 opracowano w 1990 roku, w którym ugruntowano schemat wyznaczania reszty predykcji blokowej pomiędzy klatkami, omawianej w podrozdziale~\ref{webp}, kodowane następnie, podobnie jak w przypadku JPEG, z wykorzystaniem transformacji blokowej DCT \(8 \times 8\) pikseli, kwantyzacji współczynników, zygzagowatego skanowania uzyskanych współczynników oraz kodowania entropijnego Huffmana. Dodatkowo zaimplementowano filtrowanie w pętli w celu systematycznego redukowania artefaktów blokowych. Schemat ten stał się podwaliną dla kolejnych prac grup MPEG i VCEG\cite{doi:https://doi.org/10.1002/9781118711743.ch3}.

W latach 90. wprowadzane usprawnienia do tego podejście w kodekach MPEG-1,\allowbreak MPEG-2/H.262, H.263 i MPEG-4 Part 2 --- optymalizację dla CD-ROM, obsługę wideo z przelotem, kluczową dla telewizji cyfrowej czy optymalizację dla wideokonferencji. Mały przełom nastąpił w 2003 wraz z opracowaniem formatu H.264/AVC, za sprawą bardziej rozwiniętej predykcji wewnątrzobrazowej, wprowadzenia większej elastyczności w podziale na bloki (od \(4 \times 4 \) do \(16 \times 16\) pikseli) oraz, bardzo istotne, zaawansowane kontekstowe kodowanie entropijne Context-Adaptive Variable Length Coding (CAVLC) oraz opcjonalne Context-Adaptive Binary Arithmetic Coding (CABAC). Dodatkowo struktura strumienia bitowego została zoptymalizowana dla efektywniej pakietyzacji. Wymienione techniki stanowią tylko część usprawnień względem poprzednich kodeków, a sam standard H.264 jest szeroko wykorzystywany po dziś dzień oraz doczekał się wielu wersji z kolejnymi poprawkami.

W 2013 roku grupy MPEG i VCEG wspólnie opublikował następcę standardu H.264, czyli H.265/HEVC. Miał on być odpowiedzią na rozwój technologii i idące za nią rosnące wymagania, w szczególności wzrost rozdzielczości wideo, sięgających w tym okresie już do 4K. Za kolejny wzrost efektywności wtedy najbardziej zaawansowanego z rodziny kodeka odpowiada między innymi: wielopoziomowa, drzewiasta struktura jednostek kodowania, których rozmiary znajdują się w zakresie od \(4 \times 4\) do \(64 \times 64\) pikseli, ulepszona predykcja resztkowa między klatkami, wspierająca 35 trybów kierunkowych i domyślnie stosowany CABAC. Istotną rolę pełnią także ulepszone filtry w pętli rekonstrukcji. Standard ten został zaadaptowane równie pomyślnie co H.264, znajdując wsparcie sprzętowe w większości nowoczesnych układów elektroniki użytkowej.

Technika kodowania pojedynczych obrazów (ang. intra frames) w HEVC została wykorzystana przez niezależnego inżyniera, Fabrice'a Bellarda, w 2014 roku do stworzenia formatu Better Portable Graphics (BPG). Oferuje on kompresję stratą i bezstratną, obsługę obrazów w wyższej głębi bitowej (8--14 bitów na składową). próbkowanie chrominancji w trybach 4:4:4, 4:2:2 i 4:2:0 a także obsługę kanału przezroczystości. Ponadto wspiera wiele przestrzeni barw (YCbCr z BT.601/709/2020, RGB, YCgCo, CMYK), co pozwala na wykorzystanie w HDR. Wykorzystanie tych oraz innych mechanizmów kodeka HEVC czyni ten format bardziej wydajnym i wszechstronnym nie tylko w porównaniu z dominującym wówczas, lecz archaicznym JPEG, ale także z bardziej nowoczesnym WebP. Niestety, mimo technicznej przewagi, ochrona patentowa standardu HEVC istotnie utrudniła adaptacją niezależnego rozwiązania i z przyczyn formalnych BPG nigdy nie znalazł szerokiego zastosowania w cyfrowym świecie. Obecnie traktuje się ten format głównie jako ciekawostkę techniczną i punkt odniesienia w testach jakości.

\subsection{Kompresja obrazów oparta na uczeniu maszynowym}
Podwaliną dla obecnych rozwiązań w tej dziedzinie był rozwój grupy modeli zwaną autoenkoderami. Są to sieci nienadzorowane składające się z funkcji kodującej i funkcji dekodującej, pomiędzy którymi znajduje się przestrzeń cech ukrytych. Ich celem jest nauka efektywnego zapisywania w tej przestrzeni danych wejściowych w reprezentacji o zmniejszonej wymiarowości. Ta architektura sieci neuronowej stała się podstawowym rozwiązaniem do kompresowania danych w dziedzinie głębokiego uczenia. Jednak samo ograniczenie wymiarów danych nie jest wystarczające do uzyskania wysokiego stopnia kompresji. Choć kierunek badań został uznany za obiecujący to metoda ta wypadała gorzej od ówcześnie znanych metod kompresji obrazów\cite{earlyaecompression}, podczas gdy techniki z użyciem DCT były we wczesnej fazie rozwoju.

Przez blisko dwie dekady nie nastąpiły przełomowe odkrycia w obszarze redukcji rozmiaru obrazów. Do drugiej dekady XXI w. dziedzina stopniowo ewoluowała do czasu, gdy rozwój informatyki przyniósł odwilż w uczeniu maszynowym. Odblokowanie potencjału układów graficznych (ang. graphics processing unit, GPU), gwałtowny przyrost danych w internecie, czy udoskonalenie metod optymalizacyjnych --- te i inne czynniki przyniosły się do powstania w 2012 roku sieci neuronowych o niespotykanej dotąd efektywności w zadaniach wizji komputerowej. W tym właśnie roku opracowano głęboką sieć splotową AlexNet, której sukces otworzył drogę do dalszych badań nad tą grupą modeli. Architektura CNN znajdowała kolejne zastosowania: klasyfikacja i segmentacja obrazów, detekcja obiektów, czy rozpoznawanie twarzy. Badano też możliwość użycia CNN do superrozdzielczości --- zadania zwiększania rozdzielczości obrazu. Sukces na tym polu udowodnił, że głębokie sieci neuronowe są w stanie skutecznie analizować i odtwarzać zawartość obrazu, jednocześnie poprawiając jego jakość, na podstawie wzorców z fazy treningu. Idea kompresji nie jest zgoła odmienna. Choć kompresję i superrozdzielczość można uznać za problemy o przeciwnych kierunkach przekształceń, to gdy ten drugi porównamy z dekompresją, to możemy uznać w obu przypadkach chodzi o możliwie wierne odtworzenie obrazu z reprezentacji o niższej wymiarowości, wykorzystując wiedzę zdobytą w treningu.

Pierwszą z przełomowych prac naukowych w uczonej kompresji obrazów zaprezentowano w 2016 roku. Zespół badaczy z Google, pod przewodnictwem George'a Toderici opracował rozwiązanie oparte na rekurencyjnych sieciach neuronowych, w kombinacji z warstwami splotowymi, które przewyższyło w teście jakości, mierzoną za pomocą SSIM, kodeki JPEG i WebP. Mimo że to badanie skupiło się jedynie na nisko wymiarowych obrazach \(32 \times 32 \) pikseli, to uwiarygodniło nową kategorię kodeków względem klasycznych, a scenariusz użycia nie był odrealniony, ponieważ takimi rozdzielczościami charakteryzują się miniatury. Ponadto zespołowi udało się odnieść dwa inne sukcesy. Po pierwsze kodek charakteryzował się zmienną przepływnością --- inicjalny obraz miał bardziej ograniczony rozmiar kosztem jakości, ale możliwe było dosłanie dodatkowych bitów informacji w celu jej poprawy. Po drugie --- wąskie gardło z binaryzacją i propagacją przez wartość oczekiwaną odblokowało praktyczne użycie autoenkodera w kompresji. Sieć składała się z kodera, który produkował wartości ciągłe z przedziału $[-1, 1]$. W następnym kroku binaryzator stochastycznie mapował je na wartości $\{-1, 1\}$, ale wartość oczekiwana dla każdego $x \in [-1, 1]$ była równa $x$. Dzięki temu zabiegowi pokonano problem nieróżniczkowalności dyskretyzacji. W propagacji wstecznej, dzięki zastosowaniu techinki estymatora prostego przejścia (STE, ang. straight-through estimator) krok binaryzacji ignorowano i propagowano gradient tak, jakby tego kroku nie było. Mimo tego, że w rzeczywistości pojedyncza propagacja gradientu była obarczona błędem różnicy pomiędzy wartością zdyskretyzowaną a rzeczywistą, to w skali wielu tysięcy iteracji treningowych błędy te uśredniały się do zera, dzięki wartości oczekiwanej. Architektura kodera i dekodera w zaproponowana przez zespół Toderici'ego także nie była znana wcześniej, ponieważ zostały w nich użyte warstwy ConvLSTM (convolutional long short-term memory, sieć splotowa z długą pamięcią krótkotwałą). Warstwy te pozwalają zachowywać i nadbudowywać kontekst w kolejnych iteracjach i to za ich sprawą realizuje się progresywność modelu.

W kolejnym roku ten sam zespół zaprezentował już w pełni skalowalne pod względem rozdzielczości rozwiązanie, co otworzyło drogę do kompresji obrazów pełnowymiarowych. Autorzy zrezygnowali ze sztywnych wymiarów warstw używając na etapach analizy i syntezy przekształceń splotowych. Charakteryzują się bowiem one zdolnością do zmiany wymiaru reprezentacji, zarówno w górę jak i w dół. Druga ważna zmiana, względem kodeka dla miniatur, objęła wąskie gardło. Dostrzeżono, że kodowanie wiadomości jak dotychczas nie jest wystarczająco efektywne. a odpowiedź na pytanie jak to poprawić, leżała u podstaw tej dziedziny nauki, czyli w kodowaniu entropijnym, używanym w kodekach klasycznych. W nowej odsłonie generowana przez koder dwuwymiarowa, wielokanałowa mapa cech była binaryzowana a następnie poddawana:
\begin{itemize}
	\item maskowanemu splotowi o jądrze \(7\times7\), który dla bieżącej pozycji korzysta wyłącznie z wartości po lewej i z wierszy powyżej, czyli tych zakodowanych wcześniej, tworząc lokalny kontekst do predykcji,
	\begin{sloppypar}
	\item line-LSTM czytającego wiersz po wierszu, tworząc globalny kontekst obrazu w pamięci sieci, wzmacniany z każdą linią i w każdej kolejnej iteracji sieci.
	\end{sloppypar}
\end{itemize}
Tak skonstruowany model statystyczny generuje parametry rozkładu dla każdego bitu, a te bity następnie są kodowane za pomocą kodera arytmetycznego w oparciu to te rozkłady. Takie połączenie lokalnego i globalnego kontekstu umożliwia precyzyjne modelowanie rozkładu bitów, co znacząco zwiększa efektywność kompresji w porównaniu do wcześniejszych rozwiązań.

Kolejne z fundamentalnych dla uczonej kompresji obrazów badań także przeprowadził zespół naukowców z Google, tym razem pod przewodnictwem Johannesa Ballé. Jego rozwiązanie wprowadziło nowy paradygmat tworzenia modeli dla tego zastosowania. Zasięgnięto do teorii RD (przepływności-zniekształcenia) i na jej podstawie zredefiniowano sposób trenowania sieci. Teoria RD opisuje zależność między liczbą bitów potrzebnych do zakodowania sygnału a minimalnym możliwym zniekształceniem tego sygnału po dekompresji, umożliwiając optymalizację kompromisu między kompresją a jakością. Funkcja straty, dotąd minimalizująca zniekształcenie rekonstrukcji, otrzymała nowy składnik --- przepływność (szybkość) bitową. Zgodnie z zasadami tej teorii, optymalny kodek kompresji dąży do minimalizacji zniekształcenia przy zadanym budżecie bitowym. Nową funkcję celu zdefiniowano jako ważoną sumę obu tych składników, co można wyrazić wzorem \(L=D+\lambda*R\), gdzie \(D\) to miara błędu rekonstrukcji (np. błąd średniokwadratowy), $R$ to przepływność, a \(\lambda\) to hiperparametr kontrolujący kompromis między jakością a rozmiarem danych. Dzięki takiemu podejściu osiągnięto własność, że sieć neuronowa podczas treningu uczyła się jednocześnie, w jaki sposób jak najwierniej odtwarzać obraz oraz jak tworzyć jego wewnętrzną reprezentację w taki sposób, aby minimalizować rozmiar strumienia bitowego w kodowaniu entropijnym. Osiągnięcie to pozwoliło na optymalizację obu kluczowych, choć rozbieżnych, aspektów kompresji w jednym, spójnym procesie. Zasadniczo inaczej skonstruowano też koder entropijny --- Toderici zaproponował model kontekstowy, ale oparty na sieci rekurencyjnej, co wiązało się z dużą złożonością czasową, której nie dało się zrównoleglić --- Ballé porzucił rekurencję na rzecz szybkiego i prostego modelowania niezależnych rozkładów z użyciem operacji splotowych. Pozwoliło to znacząco przyspieszyć procesy kodowania i dekodowania, przy zachowaniu dobrego modelowania statystyk cech ukrytych. Spośród nowatorskich pomysłów tego zespołu należy na koniec wspomnieć o warstwach Uogólnionej Normalizacji Dzielącej (ang. Generalized Divisive Normalization, GDN). Dotychczas stosowane aktywacje ReLU i normalizacje zastąpiono warstwami, których zachowanie ma symulować aktywacje neuronów odpowiedzialnych za  przetwarzanie bodźców wizualnych z oczu w systemie wzrokowym człowieka. W GDN aktywacja pojedynczego neuronu jest normalizowana przez aktywacje neuronów sąsiednich, co w praktyce przynosi pozytywny efekt w dekorelowaniu danych, a co za tym idzie, zachowaniu jakości przy ograniczeniu budżetu bitowego. Szereg usprawnień zaprezentowanych w opisywanym badaniu pozwolił na uzyskaniu wyników lepszych o średnio 2--3 dB w PSNR na zbiorze Kodak w porównaniu do JPEG.

Badania Todericiego oraz, szczególnie, Ballé, są kluczowe dla zrozumienia niniejszej pracy, ponieważ czerpie ona z nich podstawowe założenia, co zostanie opisane szczegółowo w kolejnym rozdziale. Na podstawie szkieletu uformowanego przez wymienionych naukowców, powstały liczne adaptacje, modyfikujące i ulepszające niektóre aspekty, mające część wspólną. Wśród jej składowych można wymienić architekturę, składającą się z kodera, kodeka entropijnego modelującego rozkłady elementów przestrzeni cech ukrytych w celu maksymalizacji bezstratnej kompresji i dekodera, optymalizację funkcji szybkości-zniekształcenia jako funkcji celu, niezależność modeli od rozmiaru wejścia, trenowania całych modeli, od początku do końca, w jednym przepływie, czy przybliżania kwantyzacji przez STE lub jednorodny szum addytywny. Lata kolejnych badań przyniosły usprawnienia w tworzeniu koderów, tworzących mniej skorelowaną reprezentację cech ukrytych, modeli entropijnych zdolnych do lepszego odwzorowywania zależności między elementami w celu generowania lepszych rozkładów, czy poszukiwania funkcji zniekształcenia najlepiej skorelowanych z ludzką percepcją. Ponieważ różnych koncepcji powstało zbyt wiele, aby uznać za zasadne opisać je tu wszystkie, przytoczone zostaną jeszcze te, które mają wpływ na niniejszą pracę.

Kolejne badania dużą wagę kładły na optymalizację modelu entropii. Balle w 2016 roku przyjął, że rozkłady elementów reprezentacji cech ukrytych są od siebie niezależne, co nie jest zgodne ze stanem faktycznym, a było jedynie technicznym uproszczeniem. Dlatego kolejne z ich badań stanowi krok ku poprawie efektywności tego elementu modelu. W pierwszej wersji zastosowano niezależne rozkłady Gaussowskie z hiperparametrami uczonymi w trakcie treningu. W następnym badaniu wzbogacono model entropijny o dodatkową boczną ścieżkę modelującą hiperprior. Hiperkoder tworzył pomocniczą reprezentację $z$, syntetyzując reprezentację cech ukrytych $y$, a hiperdekoder, odkodowując $z$, wyznaczał hiperparametry skali dla rozkładów elementów $y$. W ten sposób model entropijny zyskał kontekst z całego obrazu do wyznaczania predykcji na poziomie szczegółów. Czyni to już z tak zaprojektowanego modelu autoenkoder wariacyjny, a skuteczność takiego rozwiązania potwierdziły kolejne testy porównawcze z poprzednim modelem, wyprzedzając go o średnio 0,5--1 dB w PSNR. Wytrenowano również model, optymalizując względem MS-SSIM, co pozwoliło na zdeklasowanie innych rozwiązań w teście porównawczym na zbiorze Kodak z użyciem tej metryki.

Wykorzystanie hiperprior'a pozwoliło na wprowadzenie kontekstu globalnego do modelowania rozkładu latentu, zaś kolejnym usprawnieniem badaczy było pomyślne wdrożenie kontekstu lokalnego. Pierwsze modele z hiperprior'em, gdy pominiemy dostarczone przez niego informacje o rozkładach, nadal traktowały wszystkie elementy reprezentacji cech ukrytych jako niezależne, podczas gdy te w istocie są silnie skorelowane przestrzennie. Znani z poprzednich prac badacze z Google w kolejnej architekturze sieci umieścili dodatkowy komponent --- model kontekstowy --- składający się z warstw splotowych używających maski do wyłączenia elementów następujących. W ten sposób model wykorzystuje wiedzę o już zakodowanych elementach do kodowania następnych, wykorzystując korelacje pomiędzy nimi.
